% !TEX root = ../Monografia.tex
\chapter{Fundamentação Teórica e Revisão Bibliográfica}
\label{cap:revisao}

Este capítulo reúne a base teórica e os trabalhos que sustentam a pesquisa. A Seção~\ref{sec:carteiras}
apresenta o problema de seleção de carteiras e o papel que a matriz de covariância desempenha nele. A
Seção~\ref{sec:estimacao} revisa os estimadores dessa matriz, do mais simples ao mais estruturado. As
Seções~\ref{sec:redes} e~\ref{sec:vinculos} tratam das redes em finanças e da evidência de que vínculos
econômicos, revelados por texto e por notícias, propagam informação entre empresas. A
Seção~\ref{sec:texto-sigma} detalha os três trabalhos que ligam texto à matriz de covariância, que são
os antecedentes diretos desta pesquisa, e a Seção~\ref{sec:llm} discute o uso de modelos de linguagem
para extrair essa informação. A Seção~\ref{sec:inferencia} apresenta os métodos de avaliação e
inferência empregados, e a Seção~\ref{sec:lacuna} situa o trabalho em relação à literatura.

%--------------------------------------------------------------------
\section{Seleção de carteiras e o papel da matriz de covariância}
\label{sec:carteiras}

Considere $p$ ativos com vetor de retornos esperados $\mu$ e matriz de covariância $\Sigma$. Na teoria
de média-variância de \citet{markowitz1952}, uma carteira de pesos $w$ tem retorno esperado $w^\top\mu$
e variância $w^\top\Sigma w$, e o investidor escolhe, entre as carteiras com o mesmo retorno esperado,
a de menor variância. Na forma com aversão ao risco $\gamma$,
\begin{equation}
  \max_{w}\; w^\top\mu - \frac{\gamma}{2}\, w^\top\Sigma w
  \quad\text{sujeito a}\quad w^\top\mathbf 1 = 1 .
  \label{eq:markowitz}
\end{equation}
O caso extremo, em que o investidor ignora os retornos esperados, é a carteira de variância mínima
global (GMV):
\begin{equation}
  w_{\mathrm{GMV}} = \frac{\Sigma^{-1}\mathbf 1}{\mathbf 1^\top\Sigma^{-1}\mathbf 1}.
  \label{eq:gmv}
\end{equation}

A teoria supõe $\mu$ e $\Sigma$ conhecidos. Na prática, ambos são estimados com uma história curta, e o
otimizador explora os erros de estimação como se fossem oportunidades: um ativo com variância
subestimada, ou com correlações que parecem baixas por acaso, recebe peso excessivo. \citet{kan2007}
quantificam a perda de utilidade fora da amostra causada pela incerteza nos parâmetros e mostram que,
com muitos ativos, o erro na estimação de $\Sigma$ tem impacto comparável ao erro na média. Os retornos
esperados são notoriamente mais difíceis de estimar, e por isso o estimador de Bayes-Stein, que encolhe
as médias amostrais para a média da carteira de variância mínima, é uma correção usual \citep{jorion1986}.
A GMV da Equação~\eqref{eq:gmv} depende apenas de $\Sigma$ e se tornou, por isso, o laboratório padrão
para comparar estimadores da matriz de covariância: a carteira construída com o estimador mais próximo
da verdadeira $\Sigma$ tende a ter a menor variância realizada. \citet{engle2006} formalizam esse
critério econômico e mostram que comparar a variância realizada de carteiras ótimas é um modo válido de
ordenar previsões de covariância e de correlação.

O tamanho do problema fica claro na comparação de \citet{demiguel2009}. Nenhum dos catorze modelos de
otimização que avaliaram superou de forma consistente, fora da amostra, a carteira ingênua que divide o
capital igualmente entre os $p$ ativos (a carteira $1/N$), em Sharpe, equivalente-certeza ou giro. Para
que a otimização compensasse o erro de estimação, seriam necessárias janelas de estimação muito mais
longas do que as disponíveis na prática. A carteira $1/N$ passou, desde então, a ser a referência
mínima de qualquer estudo empírico.

As restrições de peso também interagem com o erro de estimação. \citet{jagannathan2003} mostram que
proibir posições vendidas na GMV equivale a encolher a covariância amostral: as restrições ativas
reduzem as covariâncias estimadas dos ativos que o otimizador queria vender. Por isso restrições
``erradas'', que não fazem parte das preferências do investidor, melhoram o desempenho fora da amostra.
Esse resultado é central para interpretar comparações de estimadores em carteiras restritas, em que boa
parte da regularização já vem das restrições.

Outras carteiras dependem de $\Sigma$ de modos diferentes. A máxima diversificação de
\citet{choueifaty2008} maximiza a razão $w^\top\sigma/\sqrt{w^\top\Sigma w}$, em que $\sigma$ é o vetor
de volatilidades individuais. A paridade de risco iguala as contribuições de cada ativo para o risco
total. Problemas de razão, como o máximo Sharpe e a máxima diversificação, podem ser reescritos como
programas convexos por uma mudança de variáveis homogênea \citep{cornuejols2007}. A carteira de mínimo
\emph{tracking error} minimiza a variância do retorno ativo em relação a um índice e depende da
covariância conjunta entre os ativos e o índice.

%--------------------------------------------------------------------
\section{Estimação de matrizes de covariância}
\label{sec:estimacao}

\subsection{Covariância amostral e o problema da dimensão}

Dada uma janela de $N$ retornos $r_1,\dots,r_N\in\mathbb R^p$, a covariância amostral
$S=\frac{1}{N-1}\sum_t (r_t-\bar r)(r_t-\bar r)^\top$ é não viesada, mas tem $p(p+1)/2$ parâmetros
distintos. Quando $N<p$, $S$ é singular; quando $N$ é maior que $p$, mas da mesma ordem, $S$ é mal
condicionada. Os autovalores amostrais são mais dispersos que os verdadeiros: os maiores são
superestimados e os menores, subestimados \citep{ledoit2004jmva}. Como a GMV usa $\Sigma^{-1}$, as
direções de menor autovalor, justamente as mais mal estimadas, recebem o maior peso. O número de
condição $\kappa=\lambda_{\max}/\lambda_{\min}$ mede essa fragilidade: quanto maior, mais a inversão
amplifica o ruído \citep{lopezdeprado2016,lopezdeprado2020}.

A teoria de matrizes aleatórias torna esse diagnóstico preciso. \citet{laloux1999} compararam o
espectro de matrizes de correlação de ações com a distribuição de Marchenko-Pastur, que descreve os
autovalores de uma matriz de correlação de séries independentes. A maior parte dos autovalores
empíricos cai dentro da faixa prevista para ruído puro, e apenas poucos, associados ao mercado e a
grandes setores, se destacam. \citet{bun2017} revisam os estimadores que ``limpam'' o espectro,
preservando os autovetores amostrais e corrigindo os autovalores, entre eles o corte de autovalores
(\emph{eigenvalue clipping}) e os estimadores rotacionalmente invariantes.

\subsection{Estrutura: correlação constante, setores e fatores}

A primeira resposta ao erro de estimação foi impor estrutura. \citet{elton1973} mostraram que modelos
que supõem correlação constante entre todos os pares, ou constante dentro de grupos de empresas, preveem
a correlação futura melhor que a correlação histórica par a par. Os modelos de fatores escrevem
\begin{equation}
  \Sigma = B\,F\,B^\top + D,
  \label{eq:fatores}
\end{equation}
em que $B$ é a matriz $p\times k$ de exposições aos $k$ fatores, $F$ é a covariância dos fatores e $D$
é diagonal, com as variâncias idiossincráticas. Com $k\ll p$, o número de parâmetros cai de ordem $p^2$
para ordem $pk$. \citet{chan1999} compararam modelos de fatores na previsão de covariâncias e na
construção de carteiras de variância mínima e encontraram que poucos fatores capturam a maior parte do
ganho possível. O modelo de fator único, com o mercado como fator, é o caso mais simples. Os fatores
podem ser observáveis, como mercado, setores, tamanho e momento, ou latentes, extraídos por componentes
principais. No segundo caso, o número de fatores pode ser escolhido pelos critérios de informação de
\citet{bai2002}.

\subsection{Shrinkage linear}

\citet{ledoit2003,ledoit2004honey,ledoit2004jmva} propuseram combinar a matriz amostral com um alvo
estruturado $T$,
\begin{equation}
  \hat\Sigma = \delta\,T + (1-\delta)\,S,\qquad 0\le\delta\le 1,
  \label{eq:shrinkage}
\end{equation}
escolhendo $\delta$ para minimizar a perda quadrática esperada
$\mathbb E\|\hat\Sigma-\Sigma\|_F^2$. O alvo $T$ tem poucos parâmetros e por isso pouca variância, mas
é viesado; $S$ não tem viés, mas tem muita variância. O $\delta$ ótimo equilibra as duas fontes de erro.
Os três alvos clássicos são a identidade escalada \citep{ledoit2004jmva}, a correlação constante
\citep{ledoit2004honey} e o fator único de mercado \citep{ledoit2003}. Para o alvo de correlação
constante, \citet{ledoit2004honey} derivam
\begin{equation}
  \delta^* = \max\!\left\{0,\ \min\!\left\{\frac{\hat\kappa}{N},\ 1\right\}\right\},
  \qquad \hat\kappa = \frac{\hat\pi - \hat\rho}{\hat\gamma},
  \label{eq:delta-lw}
\end{equation}
em que $\hat\pi$ estima a soma das variâncias assintóticas dos elementos de $S$, $\hat\rho$ estima a
covariância assintótica entre o alvo e $S$, e $\hat\gamma=\|T-S\|_F^2$ mede a distância entre os dois.
Quanto mais ruidosa a amostra e mais próximo o alvo, maior o peso do alvo.

A Equação~\eqref{eq:shrinkage} admite uma leitura de Bayes empírico: $T$ é a média de um \emph{prior}
sobre $\Sigma$ e $\delta$ é a confiança relativa nele. É essa leitura que abre espaço para alvos
construídos com informação externa aos retornos, como os da Seção~\ref{sec:texto-sigma}. Quando o alvo é
uma matriz de correlação construída de fora, ela precisa ser positiva semidefinida; se não for, pode ser
substituída pela matriz de correlação mais próxima, calculada por projeções alternadas
\citep{higham2002}.

\subsection{Outras abordagens}

O \emph{shrinkage} linear move todos os autovalores na mesma proporção. \citet{ledoit2017} propõem
encolher cada autovalor de forma diferente, mantendo os autovetores amostrais, e \citet{ledoit2020} obtêm
uma fórmula analítica para esse estimador. O \emph{shrinkage} não linear é uma das referências mais fortes
da literatura e, como os demais desta seção, não usa nenhuma informação além dos retornos. Outra família
impõe esparsidade: o \emph{graphical lasso} \citep{friedman2008} penaliza a matriz de precisão
$\Sigma^{-1}$, e o estimador POET \citep{fan2013} combina fatores latentes com um limiar na covariância
residual.

Na direção bayesiana, o modelo de \citet{black1992} combina um \emph{prior} de equilíbrio com visões do
investidor, cada uma com sua confiança, e \citet{he1999} explicam a escolha da incerteza do \emph{prior}.
O modelo foi formulado para os retornos esperados, mas \citet{qian2001} mostram como incorporar visões
sobre a própria covariância. A incerteza de uma covariância estimada segue a distribuição de Wishart
\citep{anderson2003}, o que dá uma escala natural para a confiança no \emph{prior}.

%--------------------------------------------------------------------
\section{Redes e grafos em finanças}
\label{sec:redes}

Uma linha de pesquisa trata a matriz de correlação como uma rede e filtra dela a estrutura mais
relevante. \citet{mantegna1999} transformou correlações em distâncias,
$d_{ij}=\sqrt{2(1-\rho_{ij})}$, e extraiu a árvore geradora mínima do grafo completo, que agrupa as ações
por setor sem que essa informação tenha sido usada. A paridade de risco hierárquica (HRP) de
\citet{lopezdeprado2016} leva a ideia ao otimizador: agrupa os ativos hierarquicamente pela distância de
correlação e aloca o capital por bissecção recursiva, com pesos inversamente proporcionais à variância de
cada grupo. A HRP nunca inverte $\Sigma$ e é, por isso, menos sensível ao mau condicionamento.

As redes também podem vir de fora dos retornos. \citet{lu2023} constroem redes de co-negociação, em que
duas ações estão ligadas quando são negociadas em momentos próximos, e usam os grupos dessa rede para
definir quais blocos da covariância idiossincrática são estimados. O método supera blocos definidos pela
classificação setorial em volatilidade e Sharpe. É o precedente metodológico mais próximo da ideia de usar
uma rede externa para estruturar $\Sigma$, embora sem texto.

%--------------------------------------------------------------------
\section{Vínculos econômicos, texto e notícias}
\label{sec:vinculos}

Há evidência ampla de que vínculos econômicos reais propagam informação entre empresas com atraso, o
que sugere que esses vínculos também moldam a estrutura de covariância. \citet{cohen2008} usam as
relações declaradas entre empresas e seus principais clientes e mostram que o retorno dos clientes
prevê o retorno dos fornecedores no mês seguinte, efeito que atribuem à atenção limitada dos
investidores. \citet{menzly2010} encontram previsibilidade cruzada entre indústrias ligadas nas tabelas
de insumo-produto, mais forte onde há menos cobertura de analistas e menos investidores institucionais.
\citet{lee2019} documentam o mesmo fenômeno para empresas com tecnologias próximas, medidas pela
sobreposição de patentes.

O texto permite medir vínculos que nenhuma classificação oficial registra. \citet{hoberg2016} calculam
a similaridade entre as descrições de produtos nos relatórios anuais (formulário 10-K) e definem
indústrias por texto (TNIC), que mudam ao longo do tempo e capturam a concorrência melhor que os
códigos SIC e NAICS. As notícias são outra fonte. \citet{scherbina2020} identificam vínculos entre
empresas pelas co-menções em notícias e mostram que os retornos das empresas líderes preveem os das
seguidoras, mesmo sem relação de cliente e fornecedor. \citet{guo2017} mostram que a cobertura conjunta
gera transbordamento de atenção: notícias sobre uma empresa movem a atenção dos investidores para as
empresas citadas junto, o que produz co-movimento e previsibilidade.

%--------------------------------------------------------------------
\section{Texto e matriz de covariância}
\label{sec:texto-sigma}

Três trabalhos levam a informação textual diretamente à estimação de $\Sigma$ e são os antecedentes
imediatos desta pesquisa.

\citet{tang2019} medem a coocorrência de pares de empresas em notícias e mostram que ela prevê a
correlação futura entre os retornos dos pares mesmo depois de controlar a correlação passada, ou seja,
a persistência da correlação. Incorporada à estimação da covariância, a coocorrência melhora a GMV. É a
evidência mais direta de que as notícias carregam informação sobre $\Sigma$, mas a variável é pobre:
conta apenas se duas empresas aparecem juntas, sem distinguir o tipo nem o sinal da relação, e
considera só pares de empresas negociadas.

\citet{becquin2023} usam a similaridade semântica entre empresas, obtida de \emph{embeddings} de
descrições textuais e de grafos de conhecimento, para estimar a covariância. O ponto central do trabalho
é que a similaridade não deve ser usada como estimador direto, porque é descalibrada: sua escala
depende do modelo e da dimensão do \emph{embedding}. Ela serve, em vez disso, como alvo de
\emph{shrinkage}. A matriz de similaridade $\mathrm{Sim}=EE^\top$, com os \emph{embeddings} normalizados,
é lida como uma correlação; é convertida em covariância com as variâncias amostrais, $\mathrm{SimCov} =
\mathrm{Sim}\odot\sqrt{\operatorname{diag}(S)\operatorname{diag}(S)^\top}$; tem os autovalores corrigidos
para ser positiva semidefinida; e entra na Equação~\eqref{eq:shrinkage} com $\delta$ à maneira de
Ledoit e Wolf. O método foi testado nos índices S\&P 500, Nasdaq-100 e Nikkei entre 2007 e 2023, com
três meses de estimação e um de avaliação. Os autores formulam ainda a hipótese de que o ganho de um alvo
semântico é maior quando a história de preços é pouco informativa. Esse molde técnico, com o texto como
\emph{prior} e as variâncias vindas dos preços, é o seguido neste trabalho.

\citet{lu2024} usam as redes TNIC de \citet{hoberg2016} para encolher a matriz de correlação, com a
intensidade do encolhimento escolhida por aprendizado por reforço. Em carteiras de variância mínima com
mais de 400 ações e vinte anos de teste, o método supera os \emph{shrinkages} padrão em volatilidade,
Sharpe e risco de cauda, mesmo depois de custos de transação.

A Tabela~\ref{tab:antecedentes} compara esses trabalhos entre si e com a proposta desta pesquisa.

\begin{table}[htbp]
\centering
\small
\caption{Trabalhos que usam informação externa aos retornos para estimar a covariância.}
\label{tab:antecedentes}
\begin{tabularx}{\textwidth}{>{\raggedright\arraybackslash}p{2.5cm}>{\raggedright\arraybackslash}p{2.6cm}>{\raggedright\arraybackslash}X>{\centering\arraybackslash}p{1.4cm}>{\centering\arraybackslash}p{1.6cm}}
\toprule
trabalho & fonte & informação usada & tipo e sinal & entidades não negociadas \\
\midrule
\citet{tang2019} & notícias & coocorrência de pares & não & não \\
\citet{becquin2023} & descrições, grafo de conhecimento & similaridade de \emph{embeddings}, como alvo de \emph{shrinkage} & não & não \\
\citet{lu2024} & relatórios 10-K (TNIC) & similaridade de produtos, encolhimento por aprendizado por reforço & não & não \\
\citet{lu2023} & dados de negociação & co-negociação, blocos da covariância residual & não & não \\
este trabalho & notícias lidas por LLM & relações tipadas e datadas entre agentes & sim & sim \\
\bottomrule
\end{tabularx}
\end{table}

%--------------------------------------------------------------------
\section{Modelos de linguagem e informação financeira}
\label{sec:llm}

A extração de relações a partir de texto foi, por muito tempo, uma tarefa de modelos supervisionados,
treinados para um conjunto fixo de tipos de entidade e de relação, com dados anotados à mão. Os modelos
de linguagem de grande porte ajustados para seguir instruções mudaram esse quadro: um mesmo modelo pode
receber a definição da tarefa em linguagem natural e devolver relações estruturadas, com custo baixo o
bastante para processar centenas de milhares de documentos. Isso torna possível transformar um arquivo
de notícias num grafo de conhecimento, em que os nós são agentes econômicos e as arestas são as ações
que os ligam, datadas pela publicação da notícia.

O trabalho mais próximo desta pesquisa na construção do grafo é o de \citet{li2024findkg}, que ajustaram
o modelo de pesos abertos Mistral~7B \citep{jiang2023mistral} para gerar um grafo de conhecimento dinâmico a
partir de notícias financeiras. Cada fato extraído é um evento $(s,r,o,t)$: uma entidade de origem, uma
relação entre quinze tipos fixos, uma entidade de destino e a data de publicação do artigo. As entidades
incluem empresas, governos, reguladores, países, setores e indicadores econômicos. Aplicado a cerca de 400
mil artigos do \emph{Wall Street Journal} entre 1999 e 2023, o método produziu o FinDKG, com 13.645
entidades e cerca de 144 mil fatos, usado para prever arestas futuras e para montar uma carteira temática
de inteligência artificial.

O desenho do grafo desta pesquisa segue a mesma lógica: eventos datados entre agentes econômicos, tipos de
relação, entidades que vão além das empresas negociadas e o grafo consultado em janelas que avançam no
tempo. As diferenças são quatro. O modelo usado aqui não é ajustado: o Qwen2.5-14B-Instruct
\citep{qwen2024} é guiado por instruções e por decodificação restrita a um esquema JSON. Cada relação tem
tipo livre, direção e uma força entre 0 e 1, enquanto no FinDKG o sinal está embutido no tipo e não há
intensidade. As entidades são associadas aos \emph{tickers} das empresas, o que permite cruzar o grafo com
os preços par a par. E, sobretudo, o FinDKG é usado para informar o retorno esperado, enquanto este trabalho
usa o grafo para estimar a matriz de covariância, um problema que o FinDKG não aborda.

%--------------------------------------------------------------------
\section{Avaliação de estimadores e inferência}
\label{sec:inferencia}

\subsection{Avaliação de previsões e de carteiras}

A covariância verdadeira nunca é observada, e as previsões precisam ser comparadas com uma \emph{proxy}
ruidosa, como a covariância realizada no período seguinte \citep{patton2011}. O critério econômico de
\citet{engle2006}, a variância realizada da GMV, é o critério principal deste trabalho, e o $R^2$ fora da
amostra de \citet{campbell2008} mede o ganho de uma previsão sobre uma referência. Retornos diários têm
caudas pesadas e dependência temporal, o que invalida os testes clássicos de igualdade de variâncias e de
índices de Sharpe. \citet{ledoit2008} propõem testar a diferença entre dois índices de Sharpe com
\emph{bootstrap} em blocos, e \citet{ledoit2011} aplicam a mesma ideia à razão entre as variâncias de duas
carteiras. Quando a mesma hipótese é testada mais de uma vez, a correção de Bonferroni divide o nível de
significância pelo número de tentativas.

\subsection{Inferência com dados de pares}

As regressões em que a unidade é um par de ações exigem cuidado, porque pares que compartilham uma ação
não são independentes. A regressão de \citet{fama1973} estima uma regressão em corte transversal em cada
período e usa a série temporal dos coeficientes para a inferência, o que elimina a correlação entre pares
do mesmo período; os erros-padrão da série são corrigidos para autocorrelação e heterocedasticidade por
\citet{newey1987}. O procedimento de atribuição quadrática (QAP) de \citet{krackhardt1988} testa a
associação entre redes permutando juntas as linhas e colunas de uma das matrizes, o que preserva a
estrutura de dependência entre os pares.

%--------------------------------------------------------------------
\section{Lacuna e posicionamento do trabalho}
\label{sec:lacuna}

A literatura revisada estabelece três pontos. Primeiro, a estimação de $\Sigma$ é limitada pelo erro
de estimação, e impor estrutura, por fatores, setores ou \emph{shrinkage}, ajuda
(Seção~\ref{sec:estimacao}). Segundo, vínculos econômicos entre empresas, revelados por texto e por
notícias, propagam informação e geram co-movimento (Seção~\ref{sec:vinculos}). Terceiro, essa
informação pode entrar na estimação de $\Sigma$ como alvo de \emph{shrinkage} ou como estrutura de
blocos (Seção~\ref{sec:texto-sigma}).

Os trabalhos que ligam texto a $\Sigma$, porém, usam o texto de forma limitada: contam co-menções ou
medem similaridade estática entre descrições. Nenhum deles distingue o tipo da relação, se uma empresa
fornece, concorre, processa ou adquire a outra, nem o seu sinal, e nenhum considera a exposição comum
de empresas a entidades que não são negociadas, mas que movem preços, como bancos centrais,
\emph{commodities} e países. Até onde esta revisão alcançou, não há estudo que use relações extraídas
por LLM, tipadas, com sinal e com entidades não negociadas, para estimar a matriz de covariância.

Do lado da extração, \citet{li2024findkg} mostram que um LLM pode transformar notícias financeiras
num grafo de conhecimento dinâmico com tipos de relação e entidades não negociadas, mas usam esse grafo
para prever arestas e selecionar ações por tema, e não para estimar o risco conjunto dos ativos.

Este trabalho ocupa essa lacuna. Em relação a \citet{tang2019}, acrescenta o tipo, o sinal e a
exposição comum a entidades; em relação a \citet{becquin2023}, troca a similaridade estática por
relações datadas, que formam um grafo dinâmico; em relação a \citet{li2024findkg}, leva um grafo do
mesmo tipo da previsão de retornos para a estimação da matriz de covariância. E, diferentemente dos três antecedentes, avalia o
grafo sob um protocolo pré-registrado, com um placebo que separa o efeito do conteúdo das relações do
efeito genérico de regularizar a matriz, e com correção para as tentativas repetidas
(Seção~\ref{sec:inferencia}).

\clearpage
